\documentclass[10pt,a4paper]{amsart}
\usepackage{amsmath,amssymb,amsthm,graphicx}
\usepackage{verbatim}
\newenvironment{reference}{\comment}{\endcomment}
%\newenvironment{reference}{}{}
\newenvironment{slogan}{\comment}{\endcomment}
\newenvironment{history}{\comment}{\endcomment}

% For commutative diagrams we use Xy-pic
\usepackage[all]{xy}

% We use 2cell for 2-commutative diagrams.
\xyoption{2cell}
% Theorem environments.
%
\theoremstyle{plain}
\newtheorem{theorem}[subsection]{Theorem}
\newtheorem{proposition}[subsection]{Proposition}
\newtheorem{lemma}[subsection]{Lemma}

\theoremstyle{definition}
\newtheorem{definition}[subsection]{Definition}
\newtheorem{example}[subsection]{Example}
\newtheorem{exercise}[subsection]{Exercise}
\newtheorem{situation}[subsection]{Situation}

\theoremstyle{remark}
\newtheorem{remark}[subsection]{Remark}
\newtheorem{remarks}[subsection]{Remarks}

\numberwithin{equation}{subsection}

% Macros
%
\def\lim{\mathop{\mathrm{lim}}\nolimits}
\def\colim{\mathop{\mathrm{colim}}\nolimits}
\def\Spec{\mathop{\mathrm{Spec}}}
\def\Hom{\mathop{\mathrm{Hom}}\nolimits}
\def\Ext{\mathop{\mathrm{Ext}}\nolimits}
\def\SheafHom{\mathop{\mathcal{H}\!\mathit{om}}\nolimits}
\def\SheafExt{\mathop{\mathcal{E}\!\mathit{xt}}\nolimits}
\def\Sch{\mathit{Sch}}
\def\Mor{\mathop{\mathrm{Mor}}\nolimits}
\def\Ob{\mathop{\mathrm{Ob}}\nolimits}
\def\Sh{\mathop{\mathit{Sh}}\nolimits}
\def\NL{\mathop{N\!L}\nolimits}
\def\CH{\mathop{\mathrm{CH}}\nolimits}
\def\proetale{{pro\text{-}\acute{e}tale}}
\def\etale{{\acute{e}tale}}
\def\QCoh{\mathit{QCoh}}
\def\Ker{\mathop{\mathrm{Ker}}}
\def\Im{\mathop{\mathrm{Im}}}
\def\Coker{\mathop{\mathrm{Coker}}}
\def\Coim{\mathop{\mathrm{Coim}}}

% Boxtimes
%
\DeclareMathSymbol{\boxtimes}{\mathbin}{AMSa}{"02}

%
% Macros for moduli stacks/spaces
%
\def\QCohstack{\mathcal{QC}\!\mathit{oh}}
\def\Cohstack{\mathcal{C}\!\mathit{oh}}
\def\Spacesstack{\mathcal{S}\!\mathit{paces}}
\def\Quotfunctor{\mathrm{Quot}}
\def\Hilbfunctor{\mathrm{Hilb}}
\def\Curvesstack{\mathcal{C}\!\mathit{urves}}
\def\Polarizedstack{\mathcal{P}\!\mathit{olarized}}
\def\Complexesstack{\mathcal{C}\!\mathit{omplexes}}
% \Pic is the operator that assigns to X its picard group, usage \Pic(X)
% \Picardstack_{X/B} denotes the Picard stack of X over B
% \Picardfunctor_{X/B} denotes the Picard functor of X over B
\def\Pic{\mathop{\mathrm{Pic}}\nolimits}
\def\Picardstack{\mathcal{P}\!\mathit{ic}}
\def\Picardfunctor{\mathrm{Pic}}
\def\Deformationcategory{\mathcal{D}\!\mathit{ef}}
\usepackage{cleveref}
\crefname{originallemma}{original lemma}{original lemmas}
\Crefname{originallemma}{Original lemma}{Original lemmas}
\makeatletter
\newlabel{cohomology-lemma-Leray}{{01F2}{source}{}{sourcepointer-1}{}}
\newlabel{cohomology-lemma-Leray@cref}{{[originallemma][0][]01F2}{[1][0][]source}{}{}{}}
\makeatother
\crefname{originalreference}{original reference}{original references}
\Crefname{originalreference}{Original reference}{Original references}
\makeatletter
\newlabel{example-squish-tangent-vector}{{0C1K}{source}{}{sourcepointer-2}{}}
\newlabel{example-squish-tangent-vector@cref}{{[originalreference][0][]0C1K}{[1][0][]source}{}{}{}}
\makeatother
\crefname{originallemma}{original lemma}{original lemmas}
\Crefname{originallemma}{Original lemma}{Original lemmas}
\makeatletter
\newlabel{varieties-lemma-proper-geometrically-reduced-global-sections}{{0BUG}{source}{}{sourcepointer-3}{}}
\newlabel{varieties-lemma-proper-geometrically-reduced-global-sections@cref}{{[originallemma][0][]0BUG}{[1][0][]source}{}{}{}}
\makeatother
\crefname{originallemma}{original lemma}{original lemmas}
\Crefname{originallemma}{Original lemma}{Original lemmas}
\makeatletter
\newlabel{varieties-lemma-surjective-pic-birational-finite}{{0C0T}{source}{}{sourcepointer-4}{}}
\newlabel{varieties-lemma-surjective-pic-birational-finite@cref}{{[originallemma][0][]0C0T}{[1][0][]source}{}{}{}}
\makeatother
\title{Picard group of a tangent-vector squishing: exact sequence with additive kernel controlled by the squished vector}
\author{The Stacks Project Authors}
\date{stacks-project:curves.tex:local-snapshot-20261012; GNU FDL 1.2 or later}
\begin{document}
\maketitle
\noindent\textit{Editorial note.} All delivered mathematical text is the author's own, copied from the pinned Stacks Project chapter file curves.tex; no statement or proof step was rewritten, reordered or removed.\par\medskip
\noindent\textit{Editorial note.} The counted proof is the author's complete argument: the surjectivity input, the units sequence with the Leray inclusion, the kernel identification, and the injectivity argument are all retained verbatim.\par\medskip
\noindent\textit{Editorial note.} The label lemma-squish-tangent-vector names the counted statement, and the proof environment immediately following it in the source is delivered verbatim as its proof.\par\medskip
\noindent\textit{Editorial note.} The counted statement is carried by the author's own lemma environment in the chapter source.\par\medskip
\noindent\textit{Allowed standard prerequisites (editorial).}\begin{itemize}
\item Surjectivity for birational finite morphisms. Varieties, Lemma varieties-lemma-surjective-pic-birational-finite. Gives surjectivity of Pic(X) -> Pic(X') for the squishing morphism.
\item The units exact sequence of a squishing. Same chapter, Example example-squish-tangent-vector; Cohomology, Lemma cohomology-lemma-Leray. Supplies 0 -> O\_X\textasciicircum{}* -> f\_*O\_\{X'\}\textasciicircum{}* -> x\_*k -> 0 with the theta map, and the Leray bound H\textasciicircum{}1(X, f\_*O\_\{X'\}\textasciicircum{}*) subset H\textasciicircum{}1(X', O\_\{X'\}\textasciicircum{}*).
\item Regular functions on proper reduced schemes. Varieties, Lemma varieties-lemma-proper-geometrically-reduced-global-sections. Identifies global regular functions with k over the algebraically closed field, giving the injectivity statement.
\end{itemize}\medskip
\ifdefined\hypertarget\hypertarget{sourcepointer-1}{}\fi
\noindent\textit{Source reference.} Original lemma 01F2: sources/stacks-chow/tags.txt, source lines 1855--1855. Prerequisite pointer invoked at proof line(s) 3390; it names a standard fact of a sibling chapter that the counted argument uses, and every step of the counted proof built on it is retained verbatim in the delivered body.\par\medskip
\ifdefined\hypertarget\hypertarget{sourcepointer-2}{}\fi
\noindent\textit{Source reference.} Original reference 0C1K: sources/stacks-curves/pinned-original-source.tex, source lines 3228--3228. Prerequisite pointer invoked at proof line(s) 3367; it names a standard fact of the same chapter that the counted argument uses, and every step of the counted proof built on it is retained verbatim in the delivered body.\par\medskip
\ifdefined\hypertarget\hypertarget{sourcepointer-3}{}\fi
\noindent\textit{Source reference.} Original lemma 0BUG: sources/stacks-chow/tags.txt, source lines 14874--14874. Prerequisite pointer invoked at proof line(s) 3396; it names a standard fact of a sibling chapter that the counted argument uses, and every step of the counted proof built on it is retained verbatim in the delivered body.\par\medskip
\ifdefined\hypertarget\hypertarget{sourcepointer-4}{}\fi
\noindent\textit{Source reference.} Original lemma 0C0T: sources/stacks-chow/tags.txt, source lines 15098--15098. Prerequisite pointer invoked at proof line(s) 3377; it names a standard fact of a sibling chapter that the counted argument uses, and every step of the counted proof built on it is retained verbatim in the delivered body.\par\medskip
\begin{lemma}
\label{lemma-squish-tangent-vector}
Let $k$ be an algebraically closed field. If $f : X' \to X$ is the
squishing of a tangent vector $\vartheta$ as in
Example \ref{example-squish-tangent-vector}, then
there is an exact sequence
$$
(k, +) \to \Pic(X) \to \Pic(X') \to 0
$$
and the first map is injective if $X'$ is proper and reduced.
\end{lemma}

\begin{proof}
The map $\Pic(X) \to \Pic(X')$ is surjective
by Varieties, Lemma \ref{varieties-lemma-surjective-pic-birational-finite}.
Using the short exact sequence
$$
0 \to \mathcal{O}_X^* \to f_*\mathcal{O}_{X'}^*
\xrightarrow{\vartheta} x_*k \to 0
$$
of Example \ref{example-squish-tangent-vector} we obtain
$$
H^0(X', \mathcal{O}_{X'}^*) \xrightarrow{\vartheta} k \to
H^1(X, \mathcal{O}_X^*) \to H^1(X, f_*\mathcal{O}_{X'}^*)
$$
We have $H^1(X, f_*\mathcal{O}_{X'}^*) \subset H^1(X', \mathcal{O}_{X'}^*)$
(for example by the Leray spectral sequence, see
Cohomology, Lemma \ref{cohomology-lemma-Leray}).
Hence the kernel of $\Pic(X) \to \Pic(X')$ is the
cokernel of the map $\vartheta : H^0(X', \mathcal{O}_{X'}^*) \to k$.
Because $k$ is algebraically closed any regular function on a
reduced connected proper scheme over $k$ comes from an element of $k$, see
Varieties, Lemma
\ref{varieties-lemma-proper-geometrically-reduced-global-sections}.
Thus the final statement of the lemma.
\end{proof}

\end{document}
